% Options for packages loaded elsewhere
\PassOptionsToPackage{unicode}{hyperref}
\PassOptionsToPackage{hyphens}{url}
\documentclass[
  10pt,
  a4paper,
]{article}
\usepackage{xcolor}
\usepackage[margin=22mm]{geometry}
\usepackage{amsmath,amssymb}
\setcounter{secnumdepth}{-\maxdimen} % remove section numbering
\usepackage{iftex}
\ifPDFTeX
  \usepackage[T1]{fontenc}
  \usepackage[utf8]{inputenc}
  \usepackage{textcomp} % provide euro and other symbols
\else % if luatex or xetex
  \usepackage{unicode-math} % this also loads fontspec
  \defaultfontfeatures{Scale=MatchLowercase}
  \defaultfontfeatures[\rmfamily]{Ligatures=TeX,Scale=1}
\fi
\usepackage{lmodern}
\ifPDFTeX\else
  % xetex/luatex font selection
\fi
% Use upquote if available, for straight quotes in verbatim environments
\IfFileExists{upquote.sty}{\usepackage{upquote}}{}
\IfFileExists{microtype.sty}{% use microtype if available
  \usepackage[]{microtype}
  \UseMicrotypeSet[protrusion]{basicmath} % disable protrusion for tt fonts
}{}
\makeatletter
\@ifundefined{KOMAClassName}{% if non-KOMA class
  \IfFileExists{parskip.sty}{%
    \usepackage{parskip}
  }{% else
    \setlength{\parindent}{0pt}
    \setlength{\parskip}{6pt plus 2pt minus 1pt}}
}{% if KOMA class
  \KOMAoptions{parskip=half}}
\makeatother
\setlength{\emergencystretch}{3em} % prevent overfull lines
\providecommand{\tightlist}{%
  \setlength{\itemsep}{0pt}\setlength{\parskip}{0pt}}
\usepackage{mathrsfs}
\usepackage{mathtools}
\usepackage{xurl}
\emergencystretch=3em
\setcounter{tocdepth}{1}
\newcommand{\Beta}{\mathrm{B}}
\DeclareUnicodeCharacter{00B2}{\ensuremath{^2}}
\DeclareUnicodeCharacter{00B1}{\ensuremath{\pm}}
\DeclareUnicodeCharacter{00D7}{\ensuremath{\times}}
\DeclareUnicodeCharacter{2192}{\ensuremath{\to}}
\DeclareUnicodeCharacter{21A6}{\ensuremath{\mapsto}}
\DeclareUnicodeCharacter{2208}{\ensuremath{\in}}
\DeclareUnicodeCharacter{2209}{\ensuremath{\notin}}
\DeclareUnicodeCharacter{2212}{\ensuremath{-}}
\DeclareUnicodeCharacter{2227}{\ensuremath{\wedge}}
\DeclareUnicodeCharacter{2228}{\ensuremath{\vee}}
\DeclareUnicodeCharacter{2260}{\ensuremath{\ne}}
\DeclareUnicodeCharacter{2264}{\ensuremath{\le}}
\DeclareUnicodeCharacter{2265}{\ensuremath{\ge}}
\DeclareUnicodeCharacter{2297}{\ensuremath{\otimes}}
\DeclareUnicodeCharacter{00B3}{\ensuremath{^3}}
\DeclareUnicodeCharacter{2074}{\ensuremath{^4}}
\DeclareUnicodeCharacter{2076}{\ensuremath{^6}}
\DeclareUnicodeCharacter{03BB}{\ensuremath{\lambda}}
\DeclareUnicodeCharacter{03BC}{\ensuremath{\mu}}
\DeclareUnicodeCharacter{03B5}{\ensuremath{\epsilon}}
\DeclareUnicodeCharacter{03C4}{\ensuremath{\tau}}
\DeclareUnicodeCharacter{03A6}{\ensuremath{\Phi}}
\DeclareUnicodeCharacter{03B1}{\ensuremath{\alpha}}
\DeclareUnicodeCharacter{03C3}{\ensuremath{\sigma}}
\DeclareUnicodeCharacter{03B4}{\ensuremath{\delta}}
\DeclareUnicodeCharacter{03C0}{\ensuremath{\pi}}
\DeclareUnicodeCharacter{039B}{\ensuremath{\Lambda}}
\DeclareUnicodeCharacter{2202}{\ensuremath{\partial}}
\DeclareUnicodeCharacter{0393}{\ensuremath{\Gamma}}
\DeclareUnicodeCharacter{2205}{\ensuremath{\varnothing}}
\DeclareUnicodeCharacter{221E}{\ensuremath{\infty}}
\DeclareUnicodeCharacter{2200}{\ensuremath{\forall}}
\DeclareUnicodeCharacter{00B9}{\ensuremath{^1}}
\DeclareUnicodeCharacter{00B0}{\ensuremath{{}^\circ}}
\DeclareUnicodeCharacter{211D}{\ensuremath{\mathbb{R}}}
\usepackage{bookmark}
\IfFileExists{xurl.sty}{\usepackage{xurl}}{} % add URL line breaks if available
\urlstyle{same}
\hypersetup{
  hidelinks,
  pdfcreator={LaTeX via pandoc}}

\author{}
\date{}

\begin{document}

{
\setcounter{tocdepth}{3}
\tableofcontents
}
\section{The actual separated zeta sheet, its explicit formula, and its
positivity
defect}\label{the-actual-separated-zeta-sheet-its-explicit-formula-and-its-positivity-defect}

This calculation uses the quotient-size interpolation in the supplied
original LaTeX. It retains the supported arithmetic zero \(e\), the
unsupported element \(\tau\), the distinct residue quotients, the
completed-function endpoint residues, every gamma pole crossed by the
contour, and the functional-equation defect. The final countertest
concerns this particular shifted zeta sheet. It is not a counterexample
to the classical Riemann hypothesis.

Source coverage for this derivation:
\texttt{globalization\ nte/\allowbreak{}2/\allowbreak{}split\_\allowbreak{}zero\_\allowbreak{}secondary\_\allowbreak{}note.\allowbreak{}tex}, lines
1111--1555, read in full;
\texttt{globalization\ nte/\allowbreak{}4/\allowbreak{}split\_\allowbreak{}zero\_\allowbreak{}defect\_\allowbreak{}preprint\ (1).\allowbreak{}tex},
lines 500--854, read in full; the standing definitions, quotient sheets,
and shifted-deformation sections of
\texttt{globalization\ nte/\allowbreak{}5/\allowbreak{}some\ considerrations.\allowbreak{}tex} were read in the
returned source text. The whole-file read of the latter was truncated in
an unrelated middle section and is not recorded as complete. The
relevant named results are \texttt{thm:\allowbreak{}universal-\allowbreak{}taylor-\allowbreak{}deformation},
\texttt{cor:\allowbreak{}jets}, \texttt{thm:\allowbreak{}arithmetic-\allowbreak{}kernels-\allowbreak{}riemann},
\texttt{cor:\allowbreak{}first-\allowbreak{}log-\allowbreak{}jet-\allowbreak{}explicit},
\texttt{prop:\allowbreak{}hurwitz-\allowbreak{}standard-\allowbreak{}package}, and the defect preprint's
``Endpoint residues'' and ``Invariant/anti-invariant decomposition.''
The secondary note lists The Clankers;
\texttt{some\ considerrations.\allowbreak{}tex} lists ``The Clankers / prepared in
continuation''; the defect preprint has a blank author field. The
source-reading ledger maintained by the root task records canonical
source IDs. This file supplies the additional proofs below; it does not
claim those older formulas as new.

\subsection{SW1. The typed quotient maps and the interpolation they
actually
give}\label{sw1.-the-typed-quotient-maps-and-the-interpolation-they-actually-give}

Let \(S=G(\mathbb Z)=\mathbb Z\sqcup\{\tau\}\), with
\(e=0_{\mathbb Z}\ne\tau\). For \(n\ge1\), the two successive
surjections are \[
S\xrightarrow{\pi_n}G(\mathbb Z/n\mathbb Z)
\xrightarrow{p_n}\mathbb Z/n\mathbb Z,
\qquad
\pi_n(\tau)=\tau_n,\quad \pi_n(a)=a\bmod n,
\quad p_n(\tau_n)=0.
\tag{SW1}
\] The first target has \(n+1\) elements, and the second has \(n\).
Their zero fibres in the source are respectively \(\{\tau\}\) and
\(\{\tau\}\cup n\mathbb Z\). These statements follow directly from the
definitions, including at \(n=1\), when the first quotient has two
elements.

The two Dirichlet sums, and the supplied interpolation between their
numerical weights, are \[
F_t(s)=\sum_{n\ge1}(n+t)^{-s}=\zeta(s,1+t),\qquad
F_0=\zeta,\quad F_1=\zeta-1,
\quad 0\le t\le1.
\tag{SW2}
\] Initially these equalities hold for \(\Re s>1\). The identity at
\(t=1\) is a change of index in an absolutely convergent series. The
intermediate values are numerical interpolation of the cardinality
weights; they are not asserted to count a nonintegral number of elements
in a quotient.

In particular, SW2 is not the separate prime-weight Euler product
\(\prod_p(1-(p+1)^{-s})^{-1}\). Its distinction from that product is
proved with the actual quotient maps in TN1--TN9 of
\href{https://github.com/KokunoYumeto/zeta-function-research-reader/blob/173dbc6ed5a03235dbe7654f928f3692107db39b/workbenches/splitzero-tandem/branches/tau-zero-prime-positivity/TAU_WEIL_NORM_RECONSTRUCTION.md}{prime-norm
comparison proof}. For example the weights satisfy
\((m+1)(n+1)-(mn+1)=m+n\); retaining the labels gives the multiplicative
vector \((1,n)\), while summing its entries gives the nonmultiplicative
number \(n+1\).

\subsection{SW2. Analytic continuation and estimates used
below}\label{sw2.-analytic-continuation-and-estimates-used-below}

Write \(a=1+t\in[1,2]\). For any integer \(M\ge1\), Euler summation with
\(2M\) derivatives gives \[
\begin{split}
F_t(s)={}&\frac{a^{1-s}}{s-1}+\frac{a^{-s}}2
 +\sum_{j=1}^{M}\frac{B_{2j}}{(2j)!}(s)_{2j-1}a^{1-s-2j}\\
&-\frac{(s)_{2M}}{(2M)!}
\int_0^\infty \widetilde B_{2M}(x)(x+a)^{-s-2M}\,dx,
\qquad \Re s>1-2M.
\end{split}
\tag{SW3}
\] Here \(\widetilde B_k(x)=B_k(\{x\})\), and
\((s)_j=s(s+1)\cdots(s+j-1)\). To check the formula, apply integration
by parts on every interval \([n,n+1]\), using \(B_k'=kB_{k-1}\),
\(B_k(1)=B_k(0)\) for \(k\ne1\), and \(B_1(1)-B_1(0)=1\). Summing the
first boundary jumps gives the values \((n+a)^{-s}\). The remaining
boundaries at the lower end are the displayed half-value and even
Bernoulli terms; the upper boundaries tend to zero for \(\Re s>1\). The
derivative of order \(j\) of \((x+a)^{-s}\) is
\((-1)^j(s)_j(x+a)^{-s-j}\), giving the sign of the remainder. This
proves SW3 first for \(\Re s>1\). Since \(\widetilde B_{2M}\) is
bounded, its integral is absolutely and locally uniformly convergent for
\(\Re s>1-2M\), so it proves the asserted continuation there.

Letting \(M\) increase proves that \(F_t\) is meromorphic on
\(\mathbb C\) with its sole pole at 1, simple of residue 1. On each
fixed vertical strip, away from that pole, SW3 bounds \(F_t(s)\) by a
fixed polynomial in \(1+|\Im s|\), uniformly for \(t\in[0,1]\). Indeed
choose \(2M\) larger than one minus the left edge; the integral of
\((x+a)^{-\Re s-2M}\) is bounded uniformly on that strip, and
\((s)_{2M}\) is a polynomial. Derivatives in \(s\) satisfy the same kind
of bound by Cauchy's formula in a slightly larger strip. These
estimates, rather than a numerical approximation to a zero sum, justify
every infinite vertical integral below.

\subsection{SW3. The first infinitesimal is explicit and does not
vanish}\label{sw3.-the-first-infinitesimal-is-explicit-and-does-not-vanish}

Termwise differentiation, dominated uniformly on compact subsets of
\(\Re s>1\), gives \[
\partial_t F_t(s)=-sF_t(s+1),\qquad
\partial_t^jF_t(s)=(-1)^j(s)_j F_t(s+j).
\tag{SW4}
\] Analytic continuation gives these meromorphic identities everywhere.
In the dual-number ring \(\mathbb C[\epsilon]/(\epsilon^2)\), the
precise specialization is \[
F_\epsilon(s)=\zeta(s)-\epsilon s\zeta(s+1),\qquad
\log F_\epsilon=\log\zeta-\epsilon Q,
\quad Q(s)=s\frac{\zeta(s+1)}{\zeta(s)}.
\tag{SW5}
\] The logarithm equality is local on any open set where
\(\zeta\ne0,\infty\); its differentiated form is a global meromorphic
identity. In particular \[
\left.\partial_t\frac{F_t'}{F_t}\right|_{t=0}=-Q'(s),
\qquad
\left.\partial_t\left(-\frac{F_t'}{F_t}\right)\right|_{t=0}=Q'(s).
\tag{SW6}
\] The nilpotence \(\epsilon^2=0\) does not make the coefficient of
\(\epsilon\) zero: at \(s=2\), that coefficient in \(F_\epsilon\) is
\(-2\zeta(3)\ne0\).

For \(\Re s>1\), multiplying the absolutely convergent series of
\(1/\zeta(s)\) and \(\zeta(s+1)\) yields \[
Q(s)=s\sum_{n\ge1}\alpha(n)n^{-s},\quad
\alpha(1)=1,\qquad
\alpha(n)=(-1)^{\omega(n)}\frac{\prod_{p\mid n}(p-1)}n.
\tag{SW7}
\] The local factor is
\((1-p^{-s})/(1-p^{-s-1})=1-\sum_{k\ge1}(p-1)p^{-k}p^{-ks}\), which
proves the displayed coefficients. Their series is absolutely convergent
on \(\Re s>1\), also after one logarithmic differentiation. Thus \[
Q'(s)=\sum_{n\ge1}\alpha(n)(1-s\log n)n^{-s}.
\tag{SW8}
\] This first jet includes all squarefree-prime support combinations in
\(\alpha(n)\); it is not a single new prime-power term.

\subsection{SW4. Endpoint residues, logarithmic orders, and the full
reflection
defect}\label{sw4.-endpoint-residues-logarithmic-orders-and-the-full-reflection-defect}

Use exactly the supplied completion \[
\Lambda_t(s)=\pi^{-s/2}\Gamma(s/2)F_t(s),\qquad
\ell_t(s)=\Lambda_t'(s)/\Lambda_t(s).
\tag{SW9}
\] SW3 at zero, or the first Bernoulli polynomial, gives
\(F_t(0)=-1/2-t\). Hence \[
\operatorname {Res}_0\Lambda_t=-1-2t,\qquad
\operatorname {Res}_1\Lambda_t=1,
\quad
\operatorname {PP}_{0,1}\Lambda_t
=\frac{1+t}{s(s-1)}-t\frac{2s-1}{s(s-1)}.
\tag{SW10}
\] Both poles are simple for \(0\le t\le1\). Consequently \[
\operatorname {Res}_0\ell_t=\operatorname {Res}_1\ell_t=-1.
\tag{SW11}
\] For proof, a nonzero principal coefficient gives
\(\Lambda_t=(s-b)^{-1}u_t(s)\), with \(u_t(b)\ne0\); its logarithmic
derivative is \(-1/(s-b)+u_t'/u_t\). Thus the weight of the endpoint in
a logarithmic explicit formula is its order, not the coefficient
\(-1-2t\) of the original completed function.

At every negative even integer \(-2j\), \(j\ge1\), the precise order is
\[
\operatorname {ord}_{-2j}\Lambda_t
=\operatorname {ord}_{-2j}F_t-1.
\tag{SW12}
\] This formula also retains exceptional parameter values at which a
Hurwitz zero cancels a gamma pole. At \(t=1\), \(F_1(-2j)=-1\), so all
these completed poles are present and simple. In particular multiplying
only by \(s(s-1)\) does not make \(\Lambda_1\) entire.

Define the additive and logarithmic reflection defects by \[
\Delta_t(s)=\Lambda_t(s)-\Lambda_t(1-s),\qquad
d_t(s)=\ell_t(s)+\ell_t(1-s)
=\frac d{ds}\log\frac{\Lambda_t(s)}{\Lambda_t(1-s)}.
\tag{SW13}
\] These are explicit meromorphic objects. Their signs follow by the
chain rule. At \(t=0\), the classical functional equation gives both
defects zero. SW5--SW6 give the complete infinitesimal logarithmic
defect \[
\dot d_0(s)=-Q'(s)-Q'(1-s).
\tag{SW14}
\] Its endpoint residues vanish because \(Q\) is regular at both
endpoints: \(Q(0)=-2\), from \(s\zeta(s+1)\to1\), and \(Q(1)=0\). This
vanishing is only a statement about these two residues. It does not set
SW14 or the moving zero divisor to zero.

\subsection{SW5. Exact contour identity before invoking any
positivity}\label{sw5.-exact-contour-identity-before-invoking-any-positivity}

For entire functions \(f,g\), set \[
f^\#(s)=\overline{f(1-\bar s)},\qquad A(s)=f^\#(s)g(s),
\quad K(u)=A(1/2+iu),\quad
k(v)=\frac1{2\pi}\int_{\mathbb R}K(u)e^{iuv}\,du.
\tag{SW15}
\] In the formulas below the entire tests have Gaussian decay, up to a
polynomial factor, in each fixed vertical strip. We construct such tests
explicitly in SW9. The more general finite-rectangle residue identity
requires no decay hypothesis.

Fix \(c>1\) and let \(a_c=1-c\). First take a finite rectangle with
these two real edges and imaginary edges \(\pm T\), avoiding
nonremovable poles on its boundary. The residue theorem gives \[
\frac1{2\pi i}\oint\ell_t(s)A(s)\,ds
=\sum_{F_t(\rho)=0\ \mathrm{inside}}m_\rho A(\rho)
-A(1)-\sum_{j\ge0:\,-2j\ \mathrm{inside}}A(-2j).
\tag{SW16}
\] Zeros and gamma poles are written separately, so exact cancellations
in SW12 occur by addition of their integer orders. All quantities in
SW16 are finite.

If \(\ell_t A\), \(\ell_t(s)A(1-s)\), and \(\ell_t(1-s)A(1-s)\) have
only finitely many nonremovable poles in the strip and their horizontal
integrals tend to zero, then taking \(T\to\infty\) gives \[
\begin{split}
Z_{t,c}(A)-A(1)-\sum_{j\ge0:\,-2j>1-c}A(-2j)
={}&\frac1{2\pi i}\int_{\Re s=c}
\ell_t(s)(A(s)+A(1-s))\,ds\\
&-\frac1{2\pi i}\int_{\Re s=c}d_t(s)A(1-s)\,ds.
\end{split}
\tag{SW17}
\] Here \(Z_{t,c}(A)\) is the finite sum of the zero terms not
annihilated by \(A\), counted with multiplicity. The actual tests in SW9
satisfy every property in this sentence by a displayed cancellation
identity, so SW17 is used below without a missing analytic assumption.
For the sign, the left upward integral becomes the right upward integral
of \(\ell_t(1-s)A(1-s)\); the positive contour subtracts it.
Substituting \(\ell_t(1-s)=d_t(s)-\ell_t(s)\) proves SW17.

\subsection{SW6. The non-Eulerian arithmetic term and the completed
formula}\label{sw6.-the-non-eulerian-arithmetic-term-and-the-completed-formula}

Put \(a=1+t\), \(r_n=(n+t)/a\) for \(n\ge2\), and choose \(c\) so large
that \(\sum_{n\ge2}r_n^{-c}<1\). Write
\(r_{\boldsymbol n}=\prod_{i=1}^k r_{n_i}\). Absolute convergence of the
logarithm series gives \[
-\frac{F_t'}{F_t}(s)=\log a+
\sum_{k\ge1}\frac{(-1)^{k+1}}k
\sum_{n_1,\ldots,n_k\ge2}
\log r_{\boldsymbol n}\,r_{\boldsymbol n}^{-s}.
\tag{SW18}
\] To verify differentiation and absolute convergence, take a slightly
smaller real part on which the sum is still less than 1. Its positive
margin bounds the differentiated geometric series. Explicitly, if
\(B(c)=\sum r_n^{-c}<1\) and \(D(c)=\sum(\log r_n)r_n^{-c}<\infty\), the
sum of the absolute differentiated terms is
\(\sum_{k\ge1}D(c)B(c)^{k-1}=D(c)/(1-B(c))\). The general-length
derivation and its signed atoms are developed independently in
\href{https://github.com/KokunoYumeto/zeta-function-research-reader/blob/173dbc6ed5a03235dbe7654f928f3692107db39b/workbenches/splitzero-tandem/branches/tau-zero-prime-positivity/NON_EULERIAN_LENGTH_DERIVATION.md}{non-Eulerian
length proof}.

Define the absolutely convergent functional \[
\begin{split}
\mathcal L_t(k)={}&2\log(a)k(0)\\
&+\sum_{k_0\ge1}\frac{(-1)^{k_0+1}}{k_0}
\sum_{n_1,\ldots,n_{k_0}\ge2}
\frac{\log r_{\boldsymbol n}}{\sqrt{r_{\boldsymbol n}}}
\left[k(\log r_{\boldsymbol n})+k(-\log r_{\boldsymbol n})\right].
\end{split}
\tag{SW19}
\] Gaussian strip decay implies \(|k(v)|\le C_R e^{-R|v|}\) for every
fixed \(R>0\): shift the vertical contour defining \(k\) by \(R\) in the
appropriate direction. Choose \(R=c-1/2\), or a slightly larger value,
and the absolute majorant just proved for SW18 proves convergence of
SW19. All exchanges of its sum and integral are therefore justified.

Let \[
w_\infty(u)=\Re\psi(1/4+iu/2)-\log\pi,\qquad
\mathcal D_{t,c}(A)=\frac1{2\pi i}\int_{\Re s=c}d_t(s)A(1-s)\,ds.
\tag{SW20}
\] Then SW17 is exactly \[
\boxed{\begin{aligned}
Z_{t,c}(A)={}&A(1)+\sum_{j\ge0:\,-2j>1-c}A(-2j)\\
&+\frac1{2\pi}\int_{\mathbb R}K(u)w_\infty(u)\,du
-\mathcal L_t(k)-\mathcal D_{t,c}(A).
\end{aligned}}
\tag{SW21}
\] For completeness,
\(\ell_t=-\tfrac12\log\pi+\tfrac12\psi(s/2)+F_t'/F_t\). The gamma part
can be moved from \(c\) to \(1/2\) without crossing a gamma pole, and
reflection of the integration variable gives \(w_\infty\). The digamma
series on \(\Re(s/2)>0\) bounds it polynomially, so Gaussian decay kills
the horizontal edges. A term \(r^{-s}\) in SW18 integrates to
\(r^{-1/2}[k(-\log r)+k(\log r)]\), by shifting its entire integrand to
the critical line. The constant gives \(2\log(a)k(0)\). Finally the
minus sign of \(\mathcal D\) is the sign derived in SW17. This proves
every sign and endpoint contribution in SW21.

At \(t=0\), \(d_0=0\), and SW18 reduces to the usual logarithmic
prime-power series. For tests annihilating the classical trivial zeros,
\(A(-2j)=0\) for \(j\ge1\). SW21 then has exactly the two endpoints,
archimedean term, and prime term in WA24. At \(t=1\), neither the
negative-even terms nor \(\mathcal D_{1,c}\) can be dropped merely
because the original endpoint principal part has been computed.

\subsection{SW7. Local variation of the original zero
packet}\label{sw7.-local-variation-of-the-original-zero-packet}

Let \(\rho\) be an actual zero of \(\zeta\), of multiplicity \(m\), and
take a small circle containing no other zero and neither endpoint. Its
continuation as a zero cluster of \(F_t\) is counted by \[
T_t(A)=\frac1{2\pi i}\oint \frac{F_t'(s)}{F_t(s)}A(s)\,ds.
\tag{SW22}
\] The compact boundary is zero-free at \(t=0\), so uniform continuity
makes it zero-free for all sufficiently small \(t\). The argument
principle proves that this is the sum of \(A\) over the moving cluster,
counted with multiplicities. Differentiating on that compact contour and
integrating an exact derivative gives \[
\dot T_0(A)=-\frac1{2\pi i}\oint Q'(s)A(s)\,ds
=\operatorname {Res}_{s=\rho}\bigl(Q(s)A'(s)\bigr).
\tag{SW23}
\] The integration-by-parts equality follows from the zero contour
integral of \((QA)'\). It does not require the zero to be simple. If it
is simple, the result is \[
\dot\rho(0)=\frac{\rho\zeta(\rho+1)}{\zeta'(\rho)},\qquad
\dot T_0(A)=\dot\rho(0)A'(\rho).
\tag{SW24}
\] For a nontrivial zero, the numerator is nonzero because \(\rho\ne0\)
and \(\Re(\rho+1)>1\). Thus a simple nontrivial zero has a nonzero first
displacement in this interpolation. For a multiple zero, SW23 retains
the full residue of order up to \(m\), rather than assuming a
differentiable labeling of its individual branches.

The arithmetic first variation on the right-hand side is explicit as
well. With the conventions of SW15, the integral of
\(-Q'(s)(A(s)+A(1-s))\) on a right line is \[
\sum_{n\ge1}\frac{\alpha(n)}{\sqrt n}
\left[
\left(\frac{\log n}2-1\right)(k(-\log n)+k(\log n))
+\log n\,(k'(-\log n)-k'(\log n))
\right].
\tag{SW25}
\] To verify it, differentiation of the Fourier integral gives
\(\frac1{2\pi}\int(1/2+iu)K(u)e^{iuv}du=\tfrac12 k(v)+k'(v)\).
Reflection replaces the second derivative term by its negative. Insert
SW7 and differentiate termwise. Absolute convergence follows from the
same exponential contour bounds for \(k,k'\). The reflection correction
SW14 remains alongside SW25; omitting it is not a deformation of the
same contour identity.

\subsection{SW8. The exact symmetric lift and the missing Hermitian
property}\label{sw8.-the-exact-symmetric-lift-and-the-missing-hermitian-property}

Let \[
G_t(s)=(s-1)F_t(s),\qquad
H_t(s)=G_t(s)G_t(1-s)=-s(s-1)F_t(s)F_t(1-s).
\tag{SW26}
\] Both \(G_t\) and \(H_t\) are entire by SW3. They have real
conjugation symmetry, and \(H_t(1-s)=H_t(s)\). The actual map is the
multiplicative norm for the involution on entire functions, \[
N:\mathcal O(\mathbb C)\longrightarrow\mathcal O(\mathbb C)^{s\mapsto1-s},
\qquad N(G)(s)=G(s)G(1-s).
\tag{SW27}
\] It satisfies \(N(G_1G_2)=N(G_1)N(G_2)\), and its divisor is the
original divisor plus its reflected divisor, with multiplicities. It is
not asserted to be additive or to preserve the zero count without
doubling. On the critical line \(H_t(s)=|G_t(s)|^2\ge0\).

The zero pairing of \(H_t\), on finite packets closed under
\(\rho\mapsto1-\bar\rho\), is Hermitian. Proof: applying complex
conjugation and this permutation to
\(\sum m_\rho\overline{f(1-\bar\rho)}g(\rho)\) exchanges \(f\) and
\(g\). By contrast, a divisor lacking that permutation symmetry need not
define a Hermitian form. The next calculation proves the failure, and a
negative direction in the symmetric lift, for the specific endpoint
\(t=1\).

\subsection{SW9. A fully explicit negative test for the separated
endpoint}\label{sw9.-a-fully-explicit-negative-test-for-the-separated-endpoint}

The exact special values are \[
F_1(-20)=-1,\qquad
F_1(-19)=\frac{174611}{6600}-1=\frac{168011}{6600}>0.
\tag{SW28}
\] They follow from \(\zeta(-n)=-B_{n+1}/(n+1)\), \(B_{21}=0\), and
\(B_{20}=-174611/330\), also obtainable from SW3. Continuity therefore
gives at least one real zero in \((-20,-19)\). Define \(\alpha\) to be
the least such real zero. This is well-defined: the endpoint values
exclude neighborhoods of the two endpoints, and a nonzero holomorphic
function has only finitely many zeros in the remaining compact interval.
Write \(m\ge1\) for its exact multiplicity, and put
\(\beta=1-\alpha\in(20,21)\).

Since \(F_1(\sigma)=\sum_{n\ge2}n^{-\sigma}>0\) for real \(\sigma>1\),
\(\beta\) is not a zero of \(F_1\). The divisor of \(H_1\), however, has
both \(\alpha\) and \(\beta\), each with multiplicity \(m\). Define \[
h(s)=(s-\alpha)^m(s-\beta)^m,\qquad
W(s)=e^{(s-1/2)^2}\frac{H_1(s)}{h(s)}.
\tag{SW29}
\] The quotient is entire, \(W(\alpha),W(\beta)\ne0\), and \(W^\#=W\).
Indeed reflection exchanges the two factors of \(h\) with total sign
\((-1)^{2m}=1\), and all the coefficients are real. SW3 proves that
\(W\) and every fixed derivative are bounded by a polynomial times
\(e^{-(\Im s)^2}\) on each fixed vertical strip; away from a bounded
region the denominator is a polynomial bounded below by a positive
constant times \(|\Im s|^{2m}\), and inside that region removability
gives boundedness.

For arbitrary prescribed \(z_\alpha,z_\beta\in\mathbb C\), take the
degree-at-most-one polynomial \[
q(s)=\frac{z_\alpha}{W(\alpha)}\frac{s-\beta}{\alpha-\beta}
+\frac{z_\beta}{W(\beta)}\frac{s-\alpha}{\beta-\alpha},
\qquad f=qW,\quad A=f^\#f.
\tag{SW30}
\] Then \(f(\alpha)=z_\alpha\) and \(f(\beta)=z_\beta\). At every other
zero of \(H_1\), \(W\) vanishes at least to the full multiplicity.
Therefore \(A\) annihilates all of those zero terms. Take the strip
\(-21<\Re s<22\). Its native and symmetric zero functionals are exactly
\[
Z_{1,22}(A)=m\overline{z_\beta}z_\alpha,
\qquad
Z_{H_1}(A)=m(\overline{z_\beta}z_\alpha+\overline{z_\alpha}z_\beta).
\tag{SW31}
\] No other zero contributes, including a multiple zero or a zero with
large imaginary part, because of the preceding vanishing orders.

For \((z_\alpha,z_\beta)=(1,-1)\), this gives \[
\boxed{Z_{1,22}(f^\#f)=-m<0,\qquad Z_{H_1}(f^\#f)=-2m<0.}
\tag{SW32}
\] For \((z_\alpha,z_\beta)=(1,i)\), it gives \[
\boxed{Z_{1,22}(f^\#f)=-im,}
\tag{SW33}
\] which is not real. Thus the original shifted divisor pairing is not
Hermitian, and the exact symmetric lift has a genuine negative
direction. Nonnegativity of \(H_1\) as a function on the critical line,
proved in SW8, does not imply nonnegativity of its zero pairing.

Here are the promised contour checks for the actual tests. Since \[
\frac{H_1'}{H_1}A
=q^\#q\,e^{2(s-1/2)^2}\frac{H_1'H_1}{h^2},
\tag{SW34}
\] the only possible nonremovable poles are at \(\alpha,\beta\). Every
other zero is canceled as an identity of meromorphic functions. SW3 and
its differentiated estimate bound the numerator polynomially on each
strip, so the horizontal edges tend to zero with Gaussian speed.
Likewise, using \(H_1=G_1G_1(1-s)\), each occurrence of \(F_1'/F_1\) or
its reflected counterpart multiplied by \(A\) cancels a factor of
\(F_1\) in \(H_1^2\). The same assertion holds for \(A(1-s)\), since its
denominator is the same reflection-invariant \(h^2\). The remaining
digamma factors have finitely many poles in this strip and at most
polynomial growth on high horizontal edges, by the digamma recurrence
and its convergent right-half-plane series. Consequently SW17--SW21
apply with no unproved horizontal zero-avoidance estimate.

At \(t=1,c=22\), the absolute logarithm expansion required in SW18
holds. For instance \[
\sum_{n\ge2}\left(\frac2{n+1}\right)^{22}
\le\left(\frac23\right)^{22}
+\int_3^\infty\left(\frac2x\right)^{22}dx
=\left(\frac23\right)^{22}\left(1+\frac3{21}\right)<1.
\tag{SW35}
\] Thus the corrected explicit formula evaluated on the negative test is
the exact equality \[
\boxed{
-m=A(1)+\sum_{j=0}^{10}A(-2j)
+\frac1{2\pi}\int_{\mathbb R}|f(1/2+iu)|^2w_\infty(u)du
-\mathcal L_1(k)-\mathcal D_{1,22}(A).
}
\tag{SW36}
\] Every term is defined in SW15 and SW19--SW20, and every series and
integral converges by the estimates above. This is a concrete positivity
analysis of the actual non-Eulerian separated sheet, including its
retained completion and reflection terms. It does not insert an
arbitrary Euler factor for the supported-zero prime.

The tested divisor is the full shifted divisor in the strip
\(-21<\Re s<22\), and its selected real zero lies outside the classical
critical strip. Thus SW32 does not establish nonpositivity for a
separately defined pairing restricted to \(0<\Re s<1\). Nor does it rule
out a different boundary subtraction or a selected positivity cone. Such
an operation must specify its map and retain the component it removes.
The exact maps for this detected component are given next.

\subsection{SW10. The retained two-point object and its positive and
negative
projections}\label{sw10.-the-retained-two-point-object-and-its-positive-and-negative-projections}

Keep the actual function \(W\) of SW29 and define the vector space \[
V=W\mathbb C[s],\qquad
E:V\longrightarrow\mathbb C^2,\quad E(f)=(f(\alpha),f(\beta)).
\tag{SW37}
\] Let \(j:\mathbb C^2\to V\) be the interpolation formula SW30. Direct
evaluation proves \(Ej=\operatorname{id}\). Since
\(W(\alpha),W(\beta)\ne0\), a polynomial multiple \(qW\) lies in
\(\ker E\) precisely when \(q\) is divisible by both distinct linear
factors. Hence \[
K:=\ker E=W(s-\alpha)(s-\beta)\mathbb C[s],\qquad
V=j(\mathbb C^2)\oplus K.
\tag{SW38}
\] The decomposition is explicit: \(f=jE(f)+(f-jE(f))\), and its second
term lies in \(K\). An element of the intersection has evaluation both
equal to its coefficient vector and zero, proving directness.

For \(z=E(f)\) and \(w=E(g)\), the native and reflected pairings factor
through these maps as \[
B_{\mathrm{nat}}(f,g)=m\overline{z_\beta}w_\alpha,\qquad
B_{\mathrm{sym}}(f,g)=m(\overline{z_\beta}w_\alpha+\overline{z_\alpha}w_\beta).
\tag{SW39}
\] The same cancellation used in SW31 proves this for arbitrary
\(f,g\in V\). Both forms vanish when either input lies in \(K\). The
descended matrix of \(B_{\mathrm{sym}}\) is
\(m\left(\begin{smallmatrix}0&1\\1&0\end{smallmatrix}\right)\), which is
invertible, so its full radical is exactly \(K\). Also
\(B_{\mathrm{sym}}=B_{\mathrm{nat}}+B_{\mathrm{nat}}^{*}\), where
\(B^*(f,g)=\overline{B(g,f)}\); this is the exact passage to the
Hermitian part, including its factor of two.

Define \[
\Pi_+=\frac12\begin{pmatrix}1&1\\1&1\end{pmatrix},\qquad
\Pi_-=\frac12\begin{pmatrix}1&-1\\-1&1\end{pmatrix},\qquad
P_\pm=j\Pi_\pm E,\quad P_K=1-jE.
\tag{SW40}
\] Multiplication gives \(\Pi_\pm^2=\Pi_\pm\), \(\Pi_+\Pi_-=0\), and
\(\Pi_++\Pi_-=1\). Using \(Ej=1\) proves that \(P_+,P_-,P_K\) are
mutually annihilating idempotents summing to \(1_V\). Their images are
respectively \(j\mathbb C(1,1)\), \(j\mathbb C(1,-1)\), and \(K\). These
three summands are orthogonal for \(B_{\mathrm{sym}}\). To verify all
signs and the orthogonality, write \[
c_+(f)=\frac{f(\alpha)+f(\beta)}2,\qquad
c_-(f)=\frac{f(\alpha)-f(\beta)}2.
\] Substitution in SW39 gives \[
B_{\mathrm{sym}}(f,g)=2m\overline{c_+(f)}c_+(g)
-2m\overline{c_-(f)}c_-(g).
\tag{SW41}
\] Thus the nonnegative cone is exactly \(\{|c_-|\le|c_+|\}\), with the
radical directions retained, and \(j\mathbb C(1,1)\oplus K\) is a
maximal nonnegative linear subspace. Indeed any vector outside it has
\(c_-\ne0\); subtracting its positive component, already in this
subspace, gives a negative vector by SW41, so it cannot be adjoined
while preserving nonnegativity.

A particular boundary subtraction can now be defined, rather than left
implicit: \[
B_+(f,g):=B_{\mathrm{sym}}(f,g)+2m\overline{c_-(f)}c_-(g)
=B_{\mathrm{sym}}(P_+f,P_+g)=2m\overline{c_+(f)}c_+(g).
\tag{SW42}
\] It is positive semidefinite, and its removed negative direction is
exactly the recorded image of \(P_-\). Equations SW37--SW42 identify the
defect object, its quotient, section, radical, orthogonal summands, and
one exact corrected form. They do not claim that this particular
subtraction is the arithmetic or prismatic choice required elsewhere in
the programme; they provide the complete finite object such a comparison
would have to map.

\subsection{SW11. Consequences for the programme's
vanishing-infinitesimal
question}\label{sw11.-consequences-for-the-programmes-vanishing-infinitesimal-question}

The exact distinctions now have proved maps. The quotient-size
construction SW1 maps to the interpolation SW2; its first-order base
change is SW5; its completed divisor is measured by SW16; its failure of
reflection is the meromorphic function SW13; and its reflected Hermitian
lift is the norm map SW27. The first-order endpoint residue variation of
the logarithmic derivative is zero, while the first-order function,
logarithmic derivative, and moving zero cluster are respectively SW5,
SW6, and SW23 and need not vanish. SW24 proves nonvanishing of the
displacement at every simple nontrivial classical zero.

The previously derived prime-weight formula TN21 concerns a different
function, a holomorphic-unit multiple of classical \(\zeta\) in
\(\Re s>0\). Its unchanged zero divisor and corrected archimedean term
cannot be transferred to \(F_1=\zeta-1\). The actual map from this
interpolation to a symmetric divisor is SW27, which adds the reflected
divisor and exposes SW32. A prismatic or deformation comparison must
carry these explicit objects and the reflection defect through its maps;
nilpotence of a parameter or vanishing of an endpoint residue does not
erase them.

\end{document}
